% !TEX root = ../aa_SoM_Chapters.tex

% ö = \%c3\%b6
% ä = \%C3\%A4

\xdef\sectiontitle{\Isectiontitle} %aiheuttama

%%%%%%%%%%%%%%%
\begin{frame}[label=1_3_a]%[label=1_3_TOC1]
\frametitle{\texorpdfstring{\culine[\sectiontitleulinecolor]{\sectiontitle}}{\sectiontitle} }

\vspace{-0.5cm}

\begin{itemize}[leftmargin=0.5cm,itemsep=2mm]
    \item[1.1] \hyperlink{1_1_a}{\IaSubsectiontitle}
    \item[1.2] \hyperlink{1_2_a}{\IbSubsectiontitle}	
    \item[1.3] \hyperlink{1_3_a}{\Alert[blue]{\IcSubsectiontitle}}	
    \item[1.4] \hyperlink{1_4_a}{\IdSubsectiontitle}
    \item[1.5] \hyperlink{1_5_a}{\IeSubsectiontitle}
    \item[1.6] \hyperlink{1_6_a}{\IfSubsectiontitle}
    \item[1.7] \hyperlink{1_7_a}{\IgSubsectiontitle}
\end{itemize}

\vspace{-1cm}
\begin{itemize}[leftmargin=9.5cm,itemsep=2mm]
    \item[1.8] \hyperlink{1_8_a}{\IhSubsectiontitle}
	\item[1.9] \hyperlink{1_9_a}{\IiSubsectiontitle}
	\item[1.10] \hyperlink{1_10_a}{\IjSubsectiontitle}
	\item[1.11] \hyperlink{1_11_a}{\IkSubsectiontitle}
\end{itemize}

\endofslide \end{frame}
%%%%%%%%%%%%%%%

%\xdef\IbSubsectiontitle{Entropy and Temperature}
\xdef\sectiontitle{\IcSubsectiontitle} 

%%%%%%%%%%%%%%%
\begin{frame}%[label=1_3_a]
\frametitle{\texorpdfstring{\culine[\sectiontitleulinecolor]{\sectiontitle}}{\sectiontitle} }

\begin{itemize}
	\item Einstein proposed a \CDAlert[violet]{simple model for solids}: 
	\item[] \begin{itemize}
	\item Atoms in a crystal can oscillate along three directions
	\item Atoms \Alert[blue]{oscillate independently} of nearby atoms (assumption/approximation) 
	\item Equilibrium state is determined by the crystalline structure 
	\item \CDAlert[red]{Oscillators are harmonic} (assumption/approximation), i.\,e. comply with \CDAlert[red]{Hooke's law}:  $F=-kx$
\end{itemize}

\end{itemize}

\xdef\loc{south}
\begin{tikzpicture}[remember picture,overlay] % anchor=south
    \node (A) [yshift=-0.25*\figYshift,xshift=-4.5cm, anchor=\loc] at (current page.\loc) {\includegraphics[page=2,height=4cm]{Figures_booklet/Tikz_main_kolumbus.pdf}}; %scale=\figscale. % 8cm max hei
\end{tikzpicture}

\xdef\loc{south}
\begin{tikzpicture}[remember picture,overlay] % anchor=south
    \node (A) [yshift=-0.25*\figYshift,xshift=1cm, anchor=\loc] at (current page.\loc) {\includegraphics[page=3,height=4cm]{Figures_booklet/Tikz_main_kolumbus.pdf}}; %scale=\figscale. % 8cm max hei
\end{tikzpicture}

\xdef\loc{south}
\begin{tikzpicture}[remember picture,overlay] % anchor=south
    \node (A) [yshift=-0.25*\figYshift,xshift=5.25cm, anchor=\loc] at (current page.\loc) {\includegraphics[page=4,height=4cm]{Figures_booklet/Tikz_main_kolumbus.pdf}}; %scale=\figscale. % 8cm max hei
\end{tikzpicture}
%\xdef\loc{south}
%\begin{tikzpicture}[remember picture,overlay] % anchor=south
%    \node (A) [yshift=-0.25*\figYshift,xshift=-4cm, anchor=\loc] at (current page.\loc) {\includegraphics[page=1,height=4cm]{Figures_booklet/SOM_9_Statistical_mechanics_figures/SOM_Figure_4_aa.png}}; %scale=\figscale. % 8cm max hei
%\end{tikzpicture}

%\begin{tikzpicture}[remember picture,overlay] % anchor=south
%    \node (A) [yshift=-0.25*\figYshift,xshift=3.25cm, anchor=\loc] at (current page.\loc) {\includegraphics[page=1,height=4cm]{Figures_booklet/SOM_9_Statistical_mechanics_figures/SOM_Figure_4_bb.png}}; %scale=\figscale. % 8cm max hei
%\end{tikzpicture}

\endofslide 
%\endofsection
\end{frame}
%%%%%%%%%%%%%%%


%%%%%%%%%%%%%%%
\begin{frame}
\frametitle{\texorpdfstring{\culine[\sectiontitleulinecolor]{\sectiontitle}}{\sectiontitle} }

\begin{itemize}
	\item A solid body where \CDAlert[blue]{$N$ atoms contain $3N$ oscillators}
\item All have the same ground state frequency which depends on the crystal structure
\item \CDAlert[red]{Treating oscillators quantum mechanically}, the allowed energy states for harmonic oscillator occur with spacing of $h f_{0}$, starting from $1 / 2 h f_{0}$ 
% 6.5 cm = midway, 10 cm = 3/4 
%\vspace{2.5mm}
\item[] \begin{minipage}{7.5cm}
	\item Energies of a single oscillator are
\[
\varepsilon_{j}=\textrm{((} j+\frac{1}{2} \textrm{))}  h f_{0}
\]
\item $j=0 \leftarrow \rightarrow$ ground state
\[
\varepsilon_{j}=\frac{1}{2} h f_{0}
\]
i.\,.e. $j=$~Energy state index
\end{minipage}
\vspace{2.5mm}
\end{itemize}

%\xdef\loc{south east}
%\begin{tikzpicture}[remember picture,overlay] % anchor=south
%    \node (A) [yshift=-0.25*\figYshift,xshift=\figXshift, anchor=\loc] at (current page.\loc) {\includegraphics[page=1,height=4cm]{Figures_booklet/SOM_9_Statistical_mechanics_figures/SOM_Figure_4_bb.png}}; %scale=\figscale. % 8cm max hei
%\end{tikzpicture}

\xdef\loc{south}
\begin{tikzpicture}[remember picture,overlay] % anchor=south
    \node (A) [yshift=-0.25*\figYshift,xshift=1cm, anchor=\loc] at (current page.\loc) {\includegraphics[page=3,height=4cm]{Figures_booklet/Tikz_main_kolumbus.pdf}}; %scale=\figscale. % 8cm max hei
\end{tikzpicture}

\xdef\loc{south}
\begin{tikzpicture}[remember picture,overlay] % anchor=south
    \node (A) [yshift=-0.25*\figYshift,xshift=5.25cm, anchor=\loc] at (current page.\loc) {\includegraphics[page=4,height=4cm]{Figures_booklet/Tikz_main_kolumbus.pdf}}; %scale=\figscale. % 8cm max hei
\end{tikzpicture}

\endofslide 
%\endofsection
\end{frame}
%%%%%%%%%%%%%%%

%%%%%%%%%%%%%%%
\begin{frame}
\frametitle{\texorpdfstring{\culine[\sectiontitleulinecolor]{\sectiontitle}}{\sectiontitle} }

\begin{itemize}
	\item Number of allowed states have in principle no upper limit 
	\item Two or more oscillators can be in the same state
	\item The harmonic oscillators in the solid body continuously exchange quanta of $h f_{0}$ with each other (e.\,g. via thermal excitations)
\end{itemize}

%\xdef\loc{south west}
%\begin{tikzpicture}[remember picture,overlay] % anchor=south
%    \node (A) [yshift=2cm,xshift=-\figXshift, anchor=\loc] at (current page.\loc) {\includegraphics[page=1,width=6.8cm]{Figures_booklet/SOM_9_Statistical_mechanics_figures/SOM_Figure_4_cc.png}}; %scale=\figscale. % 8cm max hei
%\end{tikzpicture}

%\xdef\loc{south}
%\begin{tikzpicture}[remember picture,overlay] % anchor=south
%    \node (A) [yshift=-0.25*\figYshift,xshift=1cm, anchor=\loc] at (current page.\loc) {\includegraphics[page=4,height=4cm]{Figures_booklet/Tikz_main_kolumbus.pdf}}; %scale=\figscale. % 8cm max hei
%\end{tikzpicture}


\xdef\loc{south east}
\begin{tikzpicture}[remember picture,overlay] % anchor=south
    \node (A) [yshift=-0.25*\figYshift,xshift=\figXshift, anchor=\loc] at (current page.\loc) {\includegraphics[page=1,width=6.8cm]{Figures_booklet/SOM_9_Statistical_mechanics_figures/SOM_Figure_5.png}}; %scale=\figscale. % 8cm max hei
\end{tikzpicture}

\xdef\loc{west}
\begin{tikzpicture}[remember picture,overlay] % anchor=south
    \node (A) [yshift=-2cm,xshift=-\figXshift, anchor=\loc] at (current page.\loc) {\includegraphics[page=5,width=8.2cm]{Figures_booklet/Tikz_main_kolumbus.pdf}}; %scale=\figscale. % 8cm max hei
\end{tikzpicture}

\endofslide 
%\endofsection
\end{frame}
%%%%%%%%%%%%%%%


%%%%%%%%%%%%%%%
\begin{frame}
\frametitle{\texorpdfstring{\culine[\sectiontitleulinecolor]{\sectiontitle}}{\sectiontitle} }

\begin{itemize}
	\item Return to the system of harmonic \\
		oscillators exchanging energy\\
		quanta $\Leftrightarrow$ \CDAlert[red]{Einstein's solid}

	\item Assume the system contains $N$\\
	 identical but separately identifiable\\ 
	 (=distinguishable) harmonic oscillators

\item Energy of $i$:th oscillator (for simplicity,\\
	counted from ground level...)
\[
E_{n_{i}}=n_{i} \hbar \omega_{0} \quad\textrm{((} n_{i}=0,1,2, \ldots \textrm{))} 
\]
And hereby, \CDAlert[tunired]{total energy of the system ($N$ oscillators)}:
\[
 E=\nsum_{i=1}^{N} E_{n_{i}}=\nsum_{i=1}^{N} n_{i} \hbar \omega_{0}=M \hbar \omega_{0} \quad \text { \textrm{((}here: } M=\nsum_{i=1}^{N} n_{i} \textrm{))} 
\]

\end{itemize}

\xdef\loc{north east}
\begin{tikzpicture}[remember picture,overlay] % anchor=south
    \node (A) [yshift=1*\figYshift,xshift=\figXshift, anchor=\loc] at (current page.\loc) {\includegraphics[page=1,width=6.8cm]{Figures_booklet/SOM_9_Statistical_mechanics_figures/SOM_Figure_5.png}}; %scale=\figscale. % 8cm max hei
\end{tikzpicture}

\endofslide 
%\endofsection
\end{frame}
%%%%%%%%%%%%%%%

%%%%%%%%%%%%%%%
\begin{frame}
\frametitle{\texorpdfstring{\culine[\sectiontitleulinecolor]{\sectiontitle}}{\sectiontitle} }

\begin{itemize}
	\item At thermal equilibrium, the \CDAlert[red]{oscillators} in the system have different energies, in fact, \CDAlert[red]{form an energy distribution}

\item The probability for a single oscillator to be in a given energy state $i$ is obtained from statistical analysis 

\item The probability to obtain subset $a$:
\[
P_{a}=\frac{\text { number of ways to obtaining subset a }}{\text { total number of ways of obtaining all possibities }}
\]
Here we assume that the probability for a certain subset of the system follows the number of ways $W$ to have arrangements within that subset 

\item we also assume that \CDAlert[blue]{each 'way', i.\,e. each microstate,} \CDAlert[blue]{is equally probable}
%
%\item So for the oscillators in our example case, Einstein's solid, the probability that an oscillator $i$ contains energy $E_{n_{i}}\textrm{((} =n_{i} $ quanta) will be
%\[
%P_{n_{i}}=\frac{\text { number of ways energy can be distributed with fixed } n_{i}}{\text { total number of ways energy can be distributed }}
%\[
\end{itemize}

\endofslide 
%\endofsection
\end{frame}
%%%%%%%%%%%%%%%

%%%%%%%%%%%%%%%
\begin{frame}
\frametitle{\texorpdfstring{\culine[\sectiontitleulinecolor]{\sectiontitle}}{\sectiontitle} }

\begin{itemize}

\item So for the oscillators in our example case, Einstein's solid, the probability that an oscillator $i$ contains energy $E_{n_{i}}$ $(=n_{i} $ quanta) will be
\[
P_{n_{i}}=\frac{\text { number of ways energy can be distributed with fixed } n_{i}}{\text { total number of ways energy can be distributed }}
\]

\item The above equation gives us conceptual understanding of the situation, but not yet simple tools to calculate things
\implies{ more detailed \CDAlert[red]{statistical analysis} is necessary }

\item The total number of ways how $N$ quantum numbers can form the integer $M = \sum_{i=1}^N n_i$ is given by \Alert[red]{binomial coefficient}
\[
\Omega(N, M)=\frac{(M+N-1) !}{M !(N-1) !} \equiv \textrm{((}\! \begin{array}{c}
M+N-1 \\
M
\end{array} \!\textrm{))}
\]

\end{itemize}

\endofslide 
%\endofsection
\end{frame}
%%%%%%%%%%%%%%%

%%%%%%%%%%%%%%%
\begin{frame}
\frametitle{\texorpdfstring{\culine[\sectiontitleulinecolor]{\sectiontitle}}{\sectiontitle} }

\begin{itemize}
%	\item The number of ways how $N$ quantum numbers can form the integer $M = \sum_{i=1}^N n_i$ is given by \Alert[red]{binomial coefficient}
%\[
%\Omega(N, M)=\frac{(M+N-1) !}{M !(N-1) !} \equiv \textrm{((}\! \begin{array}{c}
%M+N-1 \\
%M
%\end{array} \!\textrm{))}
%\[
\item The number of ways the energy can be distributed for a fixed $n_i$ is given by 
\[
 \textrm{((} \begin{array}{c}
 \textrm{((} M-n_{i} \textrm{))} +(N-1)-1 \\
 \textrm{((} M-n_{i} \textrm{))} 
\end{array} \textrm{))} 
\]
% 6.5 cm = midway, 10 cm = 3/4 
%\vspace{2.5mm}
\item[] \begin{minipage}{6.5cm}
in total resulting in
	\[
P_{n_{i}}= \frac{\textrm{((} \begin{array}{c}
 \textrm{((} M-n_{i} \textrm{))} +(N-1)-1 \\
 \textrm{((} M-n_{i} \textrm{))} 
\end{array} \textrm{))}  }{
 \textrm{((} \begin{array}{c}
M+N-1 \\
M
\end{array} \textrm{))} }
\]
\item Problem of above equation is that it is cumbersome to use
\implies{ motivates development of \CDAlert[red]{Boltzmann distribution} }
\end{minipage}
\vspace{2.5mm}


%\item $M$ balls, $N-1$ walls, i.e. $M+N$-1 symbols:
%In how many ways can we arrange $M$ balls and $N-1$ walls into a line?
%\[
%\Omega(N, M)=\frac{(M+N-1) !}{M !(N-1) !} \equiv \textrm{((}\! \begin{array}{c}
%M+N-1 \\
%M
%\end{array} \!\textrm{))}
%\[
\end{itemize}

\xdef\loc{south east}
\begin{tikzpicture}[remember picture,overlay] % anchor=south
    \node (A) [yshift=-0.25*\figYshift,xshift=\figXshift, anchor=\loc] at (current page.\loc) {\includegraphics[page=1,width=8.2cm]{Figures_booklet/SOM_9_Statistical_mechanics_figures/SOM_Figure_6.png}}; %scale=\figscale. % 8cm max hei
\end{tikzpicture}

%\endofslide 
\endofsection
\end{frame}
%%%%%%%%%%%%%%%
